\chapter*{Peta Ketergantungan Bab}
\addcontentsline{toc}{chapter}{Peta Ketergantungan Bab}
\markboth{Peta Ketergantungan Bab}{Peta Ketergantungan Bab}

Buku ini boleh dibaca tidak berurutan, tetapi tidak semua bab berdiri sendiri.
Tabel berikut menyebut prasyarat minimum setiap bab; semua prasyarat lain boleh
diabaikan pada pembacaan pertama.

\begin{center}
\small
\begin{tabular}{@{}r p{9.0cm} p{2.8cm}@{}}
\toprule
\textbf{Bab} & \textbf{Judul} & \textbf{Prasyarat} \\
\midrule
1 & Bahasa Matematika dan Seni Membuktikan & --- \\
2 & Himpunan, Relasi, Fungsi, dan Kardinalitas & 1 \\
3 & Aksioma Bilangan Real dan Supremum & 1, 2 \\
4 & Barisan dan Limit & 3 \\
5 & Deret Tak Hingga & 4 \\
6 & Ruang Vektor dan Subruang & 2 \\
7 & Basis, Dimensi, dan Transformasi Linear & 6 \\
8 & Sistem Persamaan Linear, Determinan, dan Invers & 6 \\
9 & Nilai Eigen dan Diagonalisasi & 7, 8 \\
10 & Ortogonalitas, SVD, dan Kuadrat Terkecil & 7, 9 \\
11 & Kekontinuan dan Kekontinuan Seragam & 3, 4 \\
12 & Turunan: Satu dan Banyak Variabel & 11, 7 \\
13 & Integral Riemann dan Teorema Dasar Kalkulus & 11 \\
14 & Deret Fungsi dan Pertukaran Limit & 5, 11, 13 \\
15 & Integral Lipat, Perubahan Variabel, dan Teorema Lapangan Vektor & 10, 12, 13 \\
16 & Persamaan Diferensial dan Sistem Dinamik & 12, 14, 15 \\
17 & Peluang dan Peluang Bersyarat & 2, 5 \\
18 & Variabel Acak dan Distribusi & 17, 13 \\
19 & Hukum Bilangan Besar dan Teorema Limit Pusat & 18, 5 \\
20 & Estimasi, Interval Kepercayaan, dan Uji Hipotesis & 19 \\
21 & Regresi dan Statistika Terapan & 10, 20 \\
22 & Aritmetika Floating-Point dan Analisis Galat & 3, 5 \\
23 & Optimasi Konveks dan Dualitas & 6, 12 \\
24 & Pemrograman Linear dan Integer & 8, 23 \\
25 & Jaringan Saraf dan Turunan Balik & 10, 12, 23 \\
26 & Generalisasi, Regularisasi, dan Pelatihan & 18, 25 \\
27 & Evaluasi Model yang Jujur & 20, 26 \\
28 & Metode Formal dengan Lean 4 & 1, 2 \\
29 & Adversarial Machine Learning: Serangan dan Pertahanan & 25, 22 \\
30 & Batas Serangan: Gradien Dikaburkan dan Evasi NIDS & 27, 29 \\
31 & Kriptografi: Bilangan Bulat Modular sampai Ancaman Kuantum & 2, 3, 28 \\
32 & Privasi Diferensial dan Fondasi Matematis QML & 10, 17, 18 \\
\bottomrule
\end{tabular}
\end{center}

\phantomsection
\pdfbookmark[0]{Daftar Isi}{toc}
\tableofcontents
\cleardoublepage
\phantomsection
\pdfbookmark[0]{Daftar Gambar}{lof}
\listoffigures
